\documentclass[conference]{IEEEtran}
\IEEEoverridecommandlockouts
% The preceding line is only needed to identify funding in the first footnote. If that is unneeded, please comment it out.
\usepackage{cite}
\usepackage{amsmath,amssymb,amsfonts}
\usepackage{algorithmic}
\usepackage{graphicx}
\usepackage{textcomp}
\usepackage{xcolor}
\usepackage{booktabs}
\def\BibTeX{{\rm B\kern-.05em{\sc i\kern-.025em b}\kern-.08em
    T\kern-.1667em\lower.7ex\hbox{E}\kern-.125emX}}
\begin{document}

\title{A Comparative Evaluation of Encoder and Decoder Language Models with Parameter-Efficient Fine-Tuning (PEFT/LoRA) 
for Multi-Class Text Emotion Recognition
}

\author{\IEEEauthorblockN{Francisco Kuo Liu}
\IEEEauthorblockA{\textit{School of Computer Engineering} \\
\textit{Instituto Tecnológico de Costa Rica}\\
Cartago, Costa Rica \\
fkuo@estudiantec.cr}
\and
\IEEEauthorblockN{Jimmy Feng Liu}
\IEEEauthorblockA{\textit{School of Computer Engineering} \\
\textit{Instituto Tecnológico de Costa Rica}\\
Cartago, Costa Rica \\
jfeng@estudiantec.cr}
\and
\IEEEauthorblockN{Danielo Wu Zhong}
\IEEEauthorblockA{\textit{School of Computer Engineering} \\
\textit{Instituto Tecnológico de Costa Rica}\\
Cartago, Costa Rica \\
dwu@estudiantec.cr}
}

\maketitle

\begin{abstract}
Text-based emotion recognition remains a challenge in Natural Language Processing (NLP) due to severe class imbalance, 
contextual ambiguity, and informal language in short digital texts. While traditional encoder-based architectures and recurrent neural 
networks have served as standard baselines, modern generative Large Language Models (LLMs) offer promising capabilities for capturing
subtle affective nuances. However, full fine-tuning of these models remains computationally prohibitive. 
In this work, we present an empirical evaluation comparing traditional encoder models (BiLSTM, BERT) against intermediate-scale 
decoder architectures (Qwen, Gemma, DeepSeek) using Parameter-Efficient Fine-Tuning (PEFT/LoRA). Experiments on the 
CARER / Kaggle Emotion dataset demonstrate the trade-offs in classification performance, measure via Macro-F1, under controlled hardware constraints. 
\end{abstract}

\begin{IEEEkeywords}
Text Emotion Recognition, Large Language Models, LoRA, PEFT, Macro-F1, Multi-class Classification.
\end{IEEEkeywords}

\section{Introduction}
Text-based emotion recognition is central to applications ranging from conversational agents to consumer feedback analysis and mental health monitoring. 
However, reliably detecting fine-grained emotional states in brief, informal digital texts remains difficult due to contextual ambiguity and severe 
class imbalance. While recurrent models like BiLSTM and encoder-based architectures like BERT have long served as primary baselines, generative 
Large Language Models (LLMs) have demonstrated strong potential in capturing complex semantic variations through advanced pre-training. Despite this 
potential, the massive parameter size of modern decoders makes full fine-tuning computationally unfeasible for resource-constrained environments, 
highlighting the need for efficient adaptation techniques.

Parameter-Efficient Fine-Tuning (PEFT) methods, particularly Low-Rank Adaptation (LoRA), address these hardware bottlenecks by freezing the 
pre-trained base model and updating only a small subset of low-rank adapter weights. Although decoders such as Qwen, Gemma, and DeepSeek excel 
in general language tasks, their performance when adapted via LoRA for specialized downstream classification under high class imbalance is not 
yet fully understood. Existing benchmarks often focus on broad sentiment analysis or general classification, relying on overall accuracy, a metric 
that frequently masks poor performance on minority emotional classes. Furthermore, direct comparisons between fully fine-tuned lightweight encoders 
and LoRA-adapted intermediate decoders under identical hardware constraints remain scarce in current literature.

To address these gaps, this work evaluates the trade-offs between legacy encoder baselines and modern intermediate-scale decoders on the highly 
imbalanced CARER emotion dataset. 

Our main contributions are as follows:
\begin{itemize}
    \item \textbf{Cross-Architecture Benchmark:} We provide a systematic empirical comparison between state-of-the-art intermediate decoders (Qwen, Gemma, DeepSeek) and established encoder baselines (BiLSTM, BERT) for multi-class emotion classification.
    \item \textbf{Efficiency \& Optimization Analysis:} We evaluate the operational trade-offs of PEFT/LoRA across decoder architectures, focusing on memory footprint, parameter count, and convergence rates under strict hardware limitations.
    \item \textbf{Imbalance-Aware Evaluation:} We analyze model resilience to severe class imbalance on the CARER dataset, prioritizing Macro-F1 metrics to ensure fair assessment across underrepresented emotion classes.
    \item \textbf{Practical Deployment Guidelines:} We outline concrete performance-versus-compute trade-offs to offer actionable decision frameworks for researchers and practitioners choosing between compact encoders and LoRA-adapted decoders.
\end{itemize}

The remainder of this survey is structured as follows: Section~\ref{sec:ejeA} reviews the evolution and current approaches in text-based emotion recognition. Section~\ref{sec:ejeB} analyzes parameter-efficient adaptation methods across Encoder and Decoder architectures. Section~\ref{sec:ejeC} examines the benchmark dataset characteristics and evaluation protocols. Section~\ref{sec:methodology} outlines our proposed experimental setup and methodology for future empirical training. Section~\ref{sec:discussion} discusses preliminary trade-offs and future research directions, and Section~\ref{sec:conclusion} concludes the paper.
\section{Domain Analysis \& Evolution of Text Emotion Recognition}
\label{sec:ejeA}
\subsection{Dominant Traditional Approaches}
Text-based emotion recognition has evolved through a rapid progression of computational methodologies, which are designed to address the challenges related to semantic ambiguity and unstructured language.
Some of these computational approaches have relied on multiple computational paradigms to interpret a text and determine which emotions are shown. These foundational techniques can be broadly categorized into keyword-based, corpus-based, ruled-based, tradicional machine learning and deep learning.
\begin{itemize}
    \item \textbf{Keyword-Based Approaches}: Identify explicit emotional terms using lexical databases (such as WordNet-Affect or SentiWordNet) and apply linguistic rules to assess sentiment intensity \cite{murthy2021review}.
    
    \item \textbf{Corpus-Based Approaches}: Induce domain-specific word-emotion lexicons using statistical methods and generative models across annotated or weakly-labeled text datasets \cite{murthy2021review}.
    
    \item \textbf{Rule-Based Approaches}: Apply expert heuristic rules, Part-of-Speech tagging, and phrasal patterns to infer emotional states based on predefined linguistic structure \cite{murthy2021review}.
    
    \item \textbf{Traditional Machine Learning Approaches}: Utilize supervised algorithms, such as, Support Vector Machines (SVM), Na\"ive Bayes, Decision Trees, and $K$-Nearest Neighbors—that rely on manual feature engineering to classify text into emotional categories \cite{murthy2021review}.
    
    \item \textbf{Deep Learning Approaches}: Employ neural network architectures, such as, Recurrent Neural Networks (RNNs/LSTMs), Convolutional Neural Networks (CNNs), and Transformer-based models (e.g., BERT), to automatically learn hierarchical representations and capture complex contextual semantics directly from raw text[1].
\end{itemize}

\subsection{The Shift to Generative Decoders}

The rise of generative Large Language Models (LLMs) has fundamentally changed how we approach text classification in NLP \cite{chakriswaran2019emotion}. 
Traditionally, affective computing tasks relied on discriminative models like BERT or GRU variants that learn a direct mapping function $f(X) \to Y$ 
from input text $X$ to a predefined set of labels $Y$ \cite{zhang2024comparative}. In contrast, autoregressive decoder architectures treat classification 
as a conditional text generation problem. They model the probability of generating a target class token $y$ given a prompt sequence $X$:

\begin{equation}
    P(y \mid X) = \prod_{i=1}^{N} P(w_i \mid w_{1}, \dots, w_{i-1}, X)
\end{equation}

\noindent This shift offers significant benefits, especially when it comes to leveraging broad contextual knowledge and picking up on subtle 
emotional tone that standard encoder representations often miss \cite{saravia2018carer}.

\subsubsection{Zero-Shot and Few-Shot Prompting Paradigms}
A key feature of generative decoders is their ability to handle downstream classification without updating model weights, 
relying instead on in-context learning (ICL) \cite{cheng2026personality}. In a zero-shot setup, the model gets an instruction 
template describing the task along with the input text, using its pre-trained knowledge to output the predicted label. 
Few-shot prompting builds on this by adding $k$ example pairs directly into the prompt context:

\begin{equation}
    \text{Prompt} = \{ (X_1, Y_1), (X_2, Y_2), \dots, (X_k, Y_k), X_{\text{test}} \}
\end{equation}

While zero-shot and few-shot setups avoid task-specific training costs, their reliability for fine-grained emotion tasks varies greatly. 
Zero-shot decoders are notoriously sensitive to how prompts are written, how context examples are ordered, and how subtle affective 
context is presented. Furthermore, without parameter adaptation for a specific domain, zero-shot LLMs often struggle with subtle emotion boundaries.

\subsubsection{Supervised Fine-Tuning versus In-Context Learning}
To address the instability of in-context learning, supervised fine-tuning (SFT) adjusts the model's output probabilities to align 
directly with task-specific labels. In multi-class emotion recognition, SFT restricts the output space to valid emotion categories 
while teaching the model to handle informal digital slang and implicit emotional cues \cite{kratzwald2018deep}.

However, choosing between prompting and full fine-tuning presents clear trade-offs:

\begin{itemize}
    \item \textbf{In-Context Learning (Zero/Few-Shot):} Keeps base weights frozen and avoids training costs, but can struggle with subtle emotional nuances.
    \item \textbf{Full Supervised Fine-Tuning (SFT):} Delivers strong task alignment across imbalanced emotion classes, but requires updating billions of parameters ($>10^9$), creating massive computational and memory bottlenecks during backpropagation.
\end{itemize}

This balance demonstrates why intermediate approaches are necessary. Generative decoders provide strong semantic understanding, 
but using them effectively for emotion classification requires parameter-efficient strategies, such as Low-Rank Adaptation (LoRA), 
that deliver SFT performance without excessive computational demands \cite{nwaiwu2025peft}.

\subsection{Established Standards and Open Challenges}
Sentiment analysis and emotion recognition overlap significantly in computational linguistics, but they present very different levels of 
complexity. Distinguishing basic sentiment classification from fine-grained emotion recognition makes it easier to pinpoint the key 
challenges facing modern NLP.

\subsubsection{Established Baselines in Binary Sentiment Analysis}
Binary sentiment classification, which sorts text into simple positive or negative categories, is a well-established task. 
Earlier studies applying fine-tuned Encoders like BERT alongside GRU variants reached near-peak performance on benchmarks such as 
SST-2 and IMDb \cite{zhang2024comparative}. Because these datasets rely on balanced classes and explicit sentiment words, 
standard supervised models handle well-structured text with high reliability.

\subsubsection{Open Challenge 1: Fine-Grained Multi-Class Emotion Mapping}
Moving beyond binary sentiment requires mapping text to specific psychological states, such as Ekman's six core emotions: 
anger, fear, joy, love, sadness, and surprise. The primary difficulty here is overlapping 
class boundaries. Distinct emotions like fear and surprise frequently share identical keywords, meaning the model must rely 
on broader context rather than individual terms to tell them apart \cite{kratzwald2018deep}. Modern decoder architectures 
capture these subtle contextual shifts well, but mapping their outputs into discrete emotion categories without misclassifications 
remains difficult \cite{cheng2026personality}.

\subsubsection{Open Challenge 2: Extreme Class Imbalance in Affective Corpora}
A major bottleneck in emotion datasets like the CARER benchmark is severe class imbalance \cite{saravia2018carer}. 
Everyday language features common emotions like joy or sadness far more often than high-intensity states like surprise or fear. 
Standard cross-entropy loss naturally focuses on majority-class errors, which leads to high overall accuracy but poor performance 
on minority categories. While sampling methods and weighted losses help traditional encoders, managing severe imbalance within 
Parameter-Efficient Fine-Tuning (PEFT) frameworks, where base weights remain frozen, is still an ongoing challenge for Macro-F1 
optimization \cite{nwaiwu2025peft}.

\subsubsection{Open Challenge 3: Ambiguity in Digital Text}
Short online texts introduce noise that degrades classification accuracy. The main causes of ambiguity include:

\begin{itemize}
    \item \textbf{Implicit Emotions and Sarcasm:} Phrases without direct emotion keywords (such as \textit{"Great, another delay."}) depend on pragmatic reasoning that basic encoders often miss.
    \item \textbf{Informal Language:} Heavy use of slang, shortcuts, irregular grammar, and repeated punctuation (like \textit{"???"}) creates out-of-vocabulary problems during tokenization.
    \item \textbf{Limited Context:} Brief posts provide very little surrounding text, making it harder to determine the exact emotion of words that change meaning based on subtext.
\end{itemize}

Balancing small, efficient architectures against the deep contextual understanding needed to resolve these ambiguities within 
strict memory limits remains a major objective for practical emotion recognition systems.

\section{Parameter-Efficient Adaptation of Encoder and Decoder Architectures}
\label{sec:ejeB}

\subsection{Full Fine-Tuning vs. PEFT Paradigms}
For a long time, Full Parameter Fine-Tuning (FFT) has been the 
predominant choice for fine-tuning models due to its excellent 
performance. Nevertheless, with the arrival of Large Language Models, 
the GPU memory requirements for FFT have increased exponentially 
\cite{liu-etal-2026-look}. 
Therefore, other approaches to solve this problem have emerged, known as Parameter-Efficient Fine-Tuning (PEFT) methods. These PEFT methods were created to reduce the number of parameters needed for fine-tuning while maintaining a performance level similar to FFT \cite{liu-etal-2026-look}.

Unlike FFT, which relies on updating the full set of parameters for fine-tuning purposes, PEFT implements a different approach based on updating only a small fraction of the model's parameters. This can be achieved by:
\begin{itemize}
    \item \textbf{Addition-Based:} Introduces lightweight, trainable adapter layers or prompt tokens into the architecture while freezing base weights \cite{liu-etal-2026-look}.
    \item \textbf{Selection-Based:} Selects and updates only a specific subset of existing parameters \cite{liu-etal-2026-look}.
    \item \textbf{Reparameterization-Based:} Exploits low-rank matrix decompositions to parameterize weight updates, as seen in LoRA \cite{liu-etal-2026-look}.
\end{itemize}
PEFT paradigms offer greater computational efficiency compared to FFTs, though it is important to recognize that they also entail certain drawbacks \cite{liu-etal-2026-look}. In terms of Instruction-Following performance and Adversarial Robustness, FFTs perform better under complex circumstances, as PEFT capabilities rely on simple tasks where a reduced parameter count still captures enough label information \cite{liu-etal-2026-look}.

Specifically, empirical results from other studies demonstrate that FFT consistently outperforms popular PEFT methods across instruction-following benchmarks on various model scales \cite{liu-etal-2026-look}:
\begin{itemize}
    \item \textbf{TinyLLaMA:} FFT achieves an average MT-Bench score of 2.68, compared to LoRA (2.05) and Prefix-Tuning (1.81) \cite{liu-etal-2026-look}.
    \item \textbf{Mistral-7B:} FFT reaches an average score of 4.96, outperforming LoRA (4.74) and Prefix-Tuning (4.67) \cite{liu-etal-2026-look}.
    \item \textbf{LLaMA2-7B:} FFT maintains its advantage with an average score of 5.26, whereas LoRA and Prefix-Tuning obtain 4.91 and 4.82, respectively \cite{liu-etal-2026-look}.
\end{itemize}

Furthermore, evaluation on the AdvGLUE benchmark reveals that while PEFT methods match or slightly edge out FFT on clean datasets, FFT exhibits superior robustness against adversarial perturbations \cite{liu-etal-2026-look}:
\begin{itemize}
    \item \textbf{RoBERTa-base (Adversarial Noise):} Under adversarial attacks, FFT retains an average robustness score of 36.58, significantly outperforming LoRA (30.89) and BitFit (28.21) \cite{liu-etal-2026-look}.
    \item \textbf{RoBERTa-large (Adversarial Noise):} FFT achieves the highest overall robustness score of 55.44, compared to LoRA's 53.96 \cite{liu-etal-2026-look}.
\end{itemize}
Considering the nature of this study, it is justifiable the usage of PEFT as a Fine-Tuning method, considering that emotion recognition is a simple task, which does not requires a the complexity that FFTs can manage.


\subsection{Mechanics of Low-Rank Adaptation (LoRA)}

\subsection{Comparative Methodologies: LoRA vs. QLoRA vs. Full Fine-Tuning}

\subsection{Theoretical Limitations and Underlying Assumptions}

\section{Benchmark Analysis: The CARER Dataset \& Evaluation Protocols}
\label{sec:ejeC}

\subsection{Corpus Characteristics and Domain Scope}

\subsection{Reported Metrics and Published Baselines}

\begin{table*}[t]
\centering
\caption{Comparative Summary of Existing Approaches in Affective Computing and PEFT Architectures}
\label{tab:literature_comparison}
\small
\begin{tabular}{p{2.5cm} c p{2.8cm} p{2.5cm} p{2.0cm} p{2.5cm} p{3.2cm}}
\toprule
\textbf{Reference} & \textbf{Year} & \textbf{Technique} & \textbf{Dataset / Environment} & \textbf{Primary Metric} & \textbf{Reported Result} & \textbf{Identified Limitation} \\
\midrule

Saravia et al. \cite{saravia2018carer} & 2018 & BiLSTM + Contextual Word Embeddings & CARER (Emotion Benchmark) & Macro-F1 & 0.647 (Benchmark Baseline) & Limited long-range semantic capacity; struggles with severe class imbalance. \\
\addlinespace

Kratzwald et al. \cite{kratzwald2018deep} & 2018 & Deep Learning (BiLSTM / CNN) & Financial \& Customer Reviews & Classification Accuracy & 82.4\% Accuracy & High dependence on task-specific fine-tuning; poor zero-shot generalization. \\
\addlinespace

Chakriswaran et al. \cite{chakriswaran2019emotion} & 2019 & Lexicon-Based \& Deep Learning Survey & Multimodal / Microblogs & Accuracy / F1-Score & Qualitative SOTA Overview & Highlights systemic issues with informal slang, sarcasm, and class skew. \\
\addlinespace

Zhang \cite{zhang2024comparative} & 2024 & BERT, RNN, GRU Full Fine-Tuning & Standard Sentiment Corpora & Accuracy / Macro-F1 & BERT outperforms RNN/GRU ($>90\%$) & Full fine-tuning requires substantial memory; limited to fixed context lengths. \\
\addlinespace

Nwaiwu \cite{nwaiwu2025peft} & 2025 & LoRA, $(IA)^3$, ReFT Evaluation & Low-Resource Text Benchmarks & Macro-F1 / Memory Footprint & LoRA achieves $\sim 98\%$ of full fine-tuning & Performance varies significantly across adapter target modules and target tasks. \\
\addlinespace

Mohammed \& Kora \cite{mohammed2026deepfake} & 2026 & LLMs + LoRA (PEFT) & Text-Based Deepfake Corpora & Detection Accuracy / F1 & High parameter efficiency ($>92\%$) & Performance drops under extreme quantization or overly restricted rank ($r$). \\
\addlinespace

Cheng et al. \cite{cheng2026personality} & 2026 & LLMs In-Context Learning / SFT & Conversational Emotion Datasets & Macro-F1 & Improved multi-turn emotion tracking & Highly sensitive to prompt context; expensive inference latency for broad LLMs. \\

\bottomrule
\end{tabular}
\end{table*}

\subsection{Known Dataset Limitations and Biases}

\section{Proposed Experimental Methodology}
\label{sec:methodology}

\subsection{Experimental Framework and Hardware Constraints}

\subsection{Model Selection and Adapter Configuration}

\subsection{Pre-Processing and Imbalance Mitigation Strategy}

\subsection{Evaluation Protocol and Statistical Validation}

\section{Discussion \& Open Research Directions}
\label{sec:discussion}

\section{Conclusion}
\label{sec:conclusion}

\bibliographystyle{IEEEtran}
\bibliography{references}

\end{document}
